\clearpage
\thispagestyle{empty}
\SourceLeaf{417}\pdfbookmark[0]{Störungen der Pallas}{pallas}
\vspace*{1.6in}
\begin{center}{\Large STÖRUNGEN DER PALLAS.\par}\end{center}
\clearpage
\fancyhead[L]{\small STÖRUNGEN DER PALLAS}
\SourceLeaf{419}\pdfbookmark[1]{I. Briefwechsel}{pallas-briefwechsel}
\begin{center}
{\large [I. BRIEFWECHSEL.]\par}
\bigskip
\pdfbookmark[2]{Gauss an Olbers, 3. Januar 1806}{pallas-brief-18060103}
\textsc{Gauss} an \textsc{Olbers}, Braunschweig, 3. Januar 1806.
\end{center}
\noindent\hspace*{\parindent}.\quad.\quad.\quad.\quad.\quad Einige zufällige Umstände waren vor einiger Zeit Ursache, dass ich verschiedene Untersuchungen über die Interpolationstheorie\HistoricalNoteEight{Vgl. Band III, S. 265.]} auszuarbeiten anfing; ich glaubte dies auf einen oder ein paar Bogen bringen zu können, nun sind es aber 4 geworden, die gedruckt wohl 8 füllen könnten. Gerade die letztere grössere Hälfte ist es, von der ich bei meiner Methode, die Perturbationen zu berechnen, einen sehr nothwendigen und vortheilhaften Gebrauch mache, daher ich wohl Gelegenheit wünschte, jene Abhandlung vorher irgend wo zum Druck zu befördern.\quad.\quad.\quad.\quad.\quad.
\par\bigskip
\begin{center}
\pdfbookmark[2]{Gauss an Olbers, 6. August 1810}{pallas-brief-18100806}
\textsc{Gauss} an \textsc{Olbers}, Göttingen, 6. August 1810.
\end{center}
\noindent\hspace*{\parindent}.\quad.\quad.\quad.\quad.\quad Ich habe angefangen, einen für die Societät bestimmten Aufsatz\HistoricalNoteDouble{Vgl. Band VI, S. 1.]} über meine die Pallas betreffenden Rechnungen zu schreiben, worin ausser den numerischen Resultaten mein Verfahren, aus 4 Oppositionen eine Planetenbahn zu bestimmen (eigentlich aus 4 Längen in der Bahn die elliptischen Elemente), sowie diejenigen Abkürzungen, deren ich mich bei der Methode der kleinsten Quadrate bediene, Ihnen vielleicht nicht ganz uninteressant sein werden. Wäre der Termin für die Pariser Preisfrage noch 4 Monate weiter entfernt, so würde ich jetzt vielleicht mich noch darum bewerben; aber so ist die Zeit zu kurz, und seit 10 Monaten fehlte es zu jeder \SourceLeaf{420}etwas weitläuftigern zusammenhängenden Arbeit mir durchaus an Lust und Muth. Aber demungeachtet denke ich nun alle Theile der Methode, nach der, wie ich glaube, die Pallasstörungen zweckmässig berechnet werden können, in etwa 4 oder 5 einzelnen Mémoires für die Societät auszuarbeiten.
\par\bigskip
\begin{center}
\pdfbookmark[2]{Gauss an Olbers, 24. Oktober 1810}{pallas-brief-18101024}
\textsc{Gauss} an \textsc{Olbers}, 24. Oktober 1810.
\end{center}
\noindent\hspace*{\parindent}.\quad.\quad.\quad.\quad.\quad Nächstens werde ich bei der Societät die, wie ich glaube, schon einmal gegen Sie erwähnte Arbeit\HistoricalNoteEight{Vgl. Band VI, S. 1.]} über die Pallas vorlesen. Sie werden dann einen umständlichen Bericht darüber in unsern Anzeigen finden. Theoretisches ist auch einiges darin: theils die Entwickelung meiner schon seit mehrern Jahren gebrauchten Methode, aus 4 Oppositionen eine elliptische Bahn zu bestimmen, gleichsam ein Supplement zum 2. Abschnitt des 2. Buchs meiner Theoria; theils die Erklärung eines praktischen Kunstgriffs, wodurch bei der Methode der kleinsten Quadrate das sonst so beschwerliche Eliminiren bedeutend abgekürzt wird.\quad.\quad.\quad.\quad.\quad.\quad Ich werde daher vielleicht bald anfangen, die Störungen eines der neuen Planeten vollständiger zu berechnen. Wüsste ich, dass das Institut die Preisfrage noch einmal prorogirte, so wäre ich nicht abgeneigt, die {\AstroSymbols ⚴} dazu zu wählen, sonst muss billig die {\AstroSymbols ⚳} den Vortritt haben, wo auch, weil bald die 8. {\AstroSymbols ☍} beobachtet wird, eine grössere Satisfaction zu erwarten ist. Bei der {\AstroSymbols ⚴} würde ich doch wohl zuerst anfangen, die Elemente als variabel anzusehen und ihre Störungen während 7 Jahren durch Quadraturen zu bestimmen, wozu ich mir Formeln entworfen habe, die mir etwas bequemer scheinen als die \textsc{Laplace}schen. Falls es dann gelingt (wie es nicht anders zu erwarten ist), die 6 bisher beobachteten Oppositionen, die sich gar nicht mehr in Eine Ellipse fügen wollen, gut zu vereinigen, so würde ich nach meiner schon vor 4 oder 5 Jahren entworfenen Methode die Störungen in der sonst üblichen Form, als periodische Störungen der Länge, der Breite und des Radius vector, oder noch besser seines Logarithmen berechnen.\quad.\quad.\quad.\quad.\quad.
\par
